\section{\MakeUppercase{Introdução}}
	% nao esquecer de corrigir as acentuacoes pois meu teclado esta em ingles
		
	\subsection{Confiabilidade na industria aeronáutica}
	A manutenção aeronáutica é extremamente importante e pode ser definida como processo que certifica o contínuo funcionamento de um sistema em sua função designada e o está projetado voltado á confiabilidade e segurança \cite[p. 35]{kinnison2013aviation}. Executando os procedimentos de forma preventiva, é possível otimizar a disponibilidade de aeronaves uma vez que a inspeção pode evitar danos causem atrasos, custos excessivos ou até catastróficos. Ela pode ser estruturada com base em um plano de manutenção centrado em confiabilidade, "\textit{Reliability-Centered Maintenance}" ou RCM. Essa diretriz foi desenvolvida na década de 70 pelo exército dos Estados Unidos baseado na aviação comercial e tem como característica priorizar análise prévia da integridade de sistemas ao invés do procedimento reativo aos danos \cite[p. 34]{kinnison2013aviation}. Dentre as suas tarefas condicionais agendadas descritas pelo RCM está o monitoramento do tamanho de trincas para definir o potencial de falha de um componente e a sua taxa de crescimento para estabelecer um intervalo de inspeção \cite[p. 52]{nowlan1978maintenance}. Na aviação moderna, o RCM é amplamente usado tanto na fase de projeto e desenvolvimento, quanto na operação \cite{ludek2005rcm}. É justamente na fase de projeto que se pode melhorar o comportamento de um componente, pois na operação o que resta a ser feito é o reparo ou a troca. Pensando nisso, primeiro é necessário identificar as causas de falha para definir os pontos a serem melhorados. A falha em um componente mecânicos pode ser definida quando este não pode mais desempenhar a função na qual foi projetado. Em geral, as falhas são causadas por concentradores de tensão, tendo como vetores a geometria, defeitos na microestrutura ou ate a corrosão \cite{findlay2002aircraft}. Enquanto o modo de falha da maioria dos componentes mecânicos é a corrosão, principal modo de falha na industria aeronáutica é a fadiga,pois a longa exposição do componente a cargas oscilantes podem ocasionar a iniciação, crescimento e ate ruptura frágil de trincas na estrutura \cite{tavares2017overview}. Um projeto com tolerância a danos serve para prever o estado critico de uma estrutura considerando um defeito.
	
	\subsection{O uso de LSP para retardar o crescimento da trinca}
	O Laser Shock Peening e uma técnica que tem como objetivo introduzir tensões residuais compressivas ao longo da espessura do material, desse modo reduzindo o $K$ na ponta da trinca e consequentemente, a sua taxa de crescimento. A aplicação desse tratamento superficial em aeronaves depende de sua eficácia no alumínio, pois ele é amplamente usado na indústria aeronáutica (referenciar a porcentagem de alumínio presente em um avião), mas seu comportamento sub fratura e fadiga é difícil de modelar matematicamente devido a sua alta ductilidade, podendo sair facilmente do escopo da mecânica linear da fratura. Em experimentos com chapas de 5mm de Ti-17 tratadas com LSP feitos por \textcite{SUN2021106081}, foi observado o fechamento de trinca e a diminuição da sua taxa de crescimento, esse estudo tem como objetivo verificar se os mesmos fenômenos ocorrem em chapas de 2mm de Al-5087.
	
	\subsection{Mecânica da fratura}
    A mecânica da fratura consolidou-se como uma disciplina essencial para mitigar os   altos custos e riscos de falhas estruturais, que se tornaram mais críticas com a complexidade tecnológica moderna. Evoluindo do trabalho inicial de Griffith com vidro em 1920 para as adaptações de Irwin para metais, foi introduzido o fator de intensidade de tensão ($K$) e a taxa de liberação de energia ($G$). Diferentemente da abordagem tradicional que compara apenas tensão e resistência, a mecânica da fratura fundamenta-se em três principais variáveis — tensão aplicada, tenacidade do material e tamanho da falha —, permitindo o uso de conceitos de tolerância ao dano para prever a vida útil de componentes com trincas subcríticas, abrangendo tanto regimes elásticos lineares quanto elasto-plásticos \cite{anderson2005fracture}. A propagação de trinca por fadiga...\cite{anderson2005fracture}. O caculo e interpretação do fator de intensidade de tensão muda quando ocorre o fechamento de trinca...\cite{KELLER2021107415}.
    
    
    
    
    
    
    

	